%
% Exponential- und Logarithmusfunktionen
%
\subsection{Exponential- und Logarithmusfunktionen}
Die Exponential- und Logarithmusfunktionen werden in der klassischen
Analysis mithilfe einer Potenzreihe oder durch die Umkehrfunktion
definiert.
Beide Konzepte stehen im algebraischen Umfeld eines Differentialkörpers
nicht zur Verfügung.
Es muss also eine alternative Methode gefunden werden, diese Funktionen
zu charakterisieren.

Die Potenzreihenmethode zur Lösung von Differentialgleichungen
hat in Kapitel~\ref{chapter:differentialgleichungen} gezeigt,
dass die Potenzreihe von $e^x$ eine Folge der Differentialgleichung
für die Exponentialgleichung ist.
Die Differentialgleichung lässt sich auch in einem Differentialkörper
formulieren, es liegt daher nahe, diese anstelle der Potenzreihe
zur Charakterisierung der Exponentialfunktion zu verwenden.

\begin{definition}[Exponential- und Logarithmusfunktion]
\label{buch:intalg:elementar:def:explog}
Sei $F$ ein Differentialkörper und $f\in F$.
Ein Element $\vartheta\in F$ heisst eine \emph{Exponentialfunktion}
von $f$, wenn $D\vartheta = \vartheta\cdot Df$.
\index{Exponentialfunktion}%
Die Exponentialfunktion wird manchmal auch $\vartheta=e^f$ oder
$\vartheta=\exp f$ geschrieben.
$\vartheta$ heisst ein \emph{Logarithmus} von $f$, wenn
$D\vartheta = Df/\vartheta$ oder $\vartheta\cdot D\vartheta = Df$.
\index{Logarithmus}%
Der Logarithmus wird manchmal auch $\vartheta=\log f$ geschrieben.
\end{definition}

Schon bei der Definition des Derivationsbegriffs haben wir darauf
hingewiesen, dass es für die Derivation keine Kettenregel gibt,
weil in einem Funktionen das Konzept der Funktionszusammensetzung fehlt.
Die Differentialgleichungen der
Definition~\ref{buch:intalg:elementar:def:explog}
sind die Kettenregeln für die Exponential- bzw.~Logarithmusfunktion.
Die Kettenregel ist in diesem Rahmen also weniger eine Eigenschaft
der Deriviation, sondern der Funktionen, mit denen wir arbeiten wollen.

%
% Eindeutigkeit
%
\subsubsection{Eindeutigkeit}
Exponential- und Logarithmusfunktion sind durch die Bedingungen der
Definition~\ref{buch:intalg:elementar:def:explog} nicht eindeutig
festgelegt.
Sind $\vartheta_1$ und $\vartheta_2$ zwei Logarithmen von $f\in F$,
dann ist
\[
D(\vartheta_1-\vartheta_2)
=
\frac{Df}{f}-\frac{Df}{f}
=
0,
\]
die Differenz ist also eine Konstante $\vartheta_1-\vartheta_2\in C$.
Eine Logarithmusfunktion ist also nur bis auf eine Konstante in $C$
bestimmt.
Die Verwendung der Notation $\log f$ suggeriert zwar eine eindeutig
bestimmte Funktion, es kann aber nur eine Menge von Funktionen
$\vartheta_1 + C$ gemeint sein.

Sind andererseits $\varphi_1$ und $\varphi_2$ Exponentialfunktionen
von $f$, dann ist
\[
D
\frac{\varphi_1}{\varphi_2}
=
\frac{D\varphi_1\cdot\varphi_2 - \varphi_1\cdot D\varphi_2}{\varphi_2^2}
=
\frac{
\varphi_1\cdot Df\cdot\varphi_2 - \varphi_1\cdot \varphi_2\cdot Df
}{\varphi_2^2}
=
0,
\]
der Quotient $\varphi_1/\varphi_2$ ist also eine Konstante.
Eine Exponentialfunktion ist also nur bis auf einen konstanten
Faktor in $C$ bestimmt.
Die Notation $e^f$ suggeriert eine eineutig bestimmte Funktion,
gemeint sein kann aber nur eine Menge $C\varphi_1$ von Funktionen in $F$.

%
% Logarithmus als Umkehrfunktion der Exponentialfunktion
%
\subsubsection{Logarithmus als Umkehrfunktion der Exponentialfunktion}
Das Konzept der Zusammensetzung von Funktionen kommt in einem
Differentialkörper nicht vor, so dass die Frage, ob der
Logarithmus die Umkehrfunktion der Exponentialfunktion ist,
nicht sinnvoll gestellt werden kann.
Der Begriff der Umkehrfunktion ist nur sinnvoll, wenn man
\index{Umkehrfunktion}%
Funktionen zusammensetzen kann.

Trotzdem kann man eine Interpretation finden, nach der der Logarithmus
eine Art Umkehrfunktion ist.
Da die Exponentialfunktion und der Logarithmus durch die
Definition über ihre differentiellen Eigenschaften nicht eindeutig 
bestimmt sind, kann auch die Umkehrbeziehung nur im Sinne einer
Differentialbeziehung gelten.

Sei jetzt $f\in F$ und $\vartheta$ eine Exponentialfunktion von $f$.
Nach Definition erfüllt sie $D\vartheta=\vartheta\cdot Df$.
Weiter sei $\varphi$ ein Logarithmus von $\vartheta$, d.~h.~nach
Definition ist $D\varphi=D\vartheta/\vartheta$.
Setzt man die definierende Eigenschaft der Exponentialfunktion
in die Ableitung des Logarithmus ein, erhält man
\[
D\varphi
=
\frac{D\vartheta}{\vartheta}
=
\frac{\vartheta\cdot Df}{\vartheta}
=
Df.
\]
Dies bedeutet, dass $\varphi$ bis auf eine additive Konstante mit
$f$ übereinstimmt.
Mit der konventionellen Schreibweise für Exponentialfunktion und
Logarithmus kann man dies auch als
\[
D\log e^f
=
\frac{De^f}{e^f}
=
\frac{e^f Df}{e^f}
=
Df
\]
schreiben.
Dies bedeutet, dass $\log e^f = f+c$ mit einer Konstante $c\in C$ ist.

Umgekehrt ist
\begin{align*}
D\frac{e^{\log f}}{f}
&=
\frac{(De^{\log f})\cdot f - e^{\log f}\cdot Df}{f^2}
\\
&=
\frac{e^{\log f}D(\log f)\cdot f - e^{\log f}\cdot Df }{f^2}
\\
&=
\frac{e^{\log f} (Df/f)\cdot f - e^{\log f}\cdot Df }{f^2}
=
0,
\end{align*}
$e^{\log f}$ ist also bis auf einen konstanten Faktor die ursprüngliche
Funktion $f$.

%
% Logarithmus und Exponentialgesetze
%
\subsubsection{Logarithmus- und Exponentialgesetze}
Für die algebraische Rechnung sind Regeln nötig, mit denen Ausdrücke
mit Logarithmen und Exponentialfunktionen umgeformt werden können.
Dazu gehören die Logarithmus- und Exponentialgesetze
\index{Logarithmusgesetz}%
\index{Exponentialgesetz}%
\[
\log ab = \log a + \log b
\qquad\text{und}\qquad
e^{a+b}=e^ae^b.
\]
Da weder Logarithmus noch Exponentialfunktion eindeutig festgelegt
sind, kann man so eine einfache Beziehung nicht mehr erwarten.

Sind $f_1,f_2\in F$ Funktionen in $F$ und $\vartheta_1$ und $\vartheta_2$
Exponentialfunktionen von $f_1$ bzw.~$f_2$, dann gilt
\[
D(\vartheta_1\vartheta_2)
=
(D\vartheta_1)\vartheta_2
+
\vartheta_1(D\vartheta_1)
=
\vartheta_1 (Df_1) \vartheta_2
+
\vartheta_1\vartheta_2 (Df_2)
=
\vartheta_1\vartheta_2 D(f_1+f_2).
\]
Das Produkt $\vartheta = \vartheta_1\vartheta_2$ ist daher wegen
\[
D\vartheta = \vartheta D(f_1+f_2)
\]
eine Exponentialfunktion von $f_1+f_2$.
In der konventionellen Notation wird dies als
\[
e^{f_1} e^{f_2}
=
e^{f_1+f_2}
\]
geschrieben.
Das Produkt der Exponentialfunktionen der Faktoren ist also eine
Exponentialfunktion der Summe.
Das Exponentialgesetz ist also in diesem erweiterten Sinn immer noch
gültig.

Seien jetzt umgekehrt $\vartheta_1$ und $\vartheta_2$ Logarithmen von
$f_1$ bzw.~$f_2$.
Dann gilt
\[
D(\vartheta_1+\vartheta_2)
=
\frac{Df_1}{f_1}
+
\frac{Df_2}{f_2}
=
\frac{Df_1\cdot f_2+f_1\cdot Df_2}{f_1f_2}
=
\frac{D(f_1f_2)}{f_1f_2}.
\]
Schreibt man $f=f_1f_2$, dann ist die Summe
$\vartheta=\vartheta_1+\vartheta_2$ wegen
\[
D\vartheta = \frac{Df}{f}
\]
ein Logarithmus des Produktes $f$.
In der konventionellen Notation bedeutet dies, dass
\[
\log (f_1f_2)
=
\log f_1 + \log f_2
\]
ist.
Die Summe der Logarithmen der Faktoren ist also ein Logarithmus
des Produktes.

